% ---------------------------------------------------------------------------
% Baseline results and today's experiment
% ELEC 872 - Lab (workshop, Oct 6)
%
% Overleaf: create a new project and paste this file (or upload it).
% Compiles as-is (only needs booktabs).
% ---------------------------------------------------------------------------
\documentclass[11pt]{article}

\usepackage[margin=1in]{geometry}
\usepackage{booktabs}
\usepackage{caption}
\captionsetup{font=small, labelfont=bf}

\title{Baseline: LBP Face Recognition under\\ Lighting Changes and Rotation}
\author{ELEC 872 --- Lab}
\date{October 6, 2026}

\begin{document}
\maketitle

% ---------------------------------------------------------------------------
\section{Baseline results}

\begin{table}[htbp]
  \centering
  \caption{Baseline rank-1 recognition rate on ORL for each test condition.}
  \label{tab:baseline}
  \begin{tabular}{lr}
    \toprule
    Test condition & Rank-1 accuracy \\
    \midrule
    Normal                          & 0.975 \\
    \midrule
    Uniform brightness              & 0.975 \\
    \midrule
    Lighting ramp $s = 0.2$         & 0.960 \\
    Lighting ramp $s = 0.4$         & 0.945 \\
    Lighting ramp $s = 0.6$         & 0.895 \\
    Lighting ramp $s = 0.8$         & \textbf{0.705} \\
    Lighting ramp $s = 1.0$         & 0.335 \\
    \midrule
    Rotation $5^\circ$              & 0.970 \\
    Rotation $10^\circ$             & 0.925 \\
    Rotation $15^\circ$             & 0.800 \\
    Rotation $20^\circ$             & \textbf{0.630} \\
    Rotation $25^\circ$             & 0.510 \\
    Rotation $30^\circ$             & 0.340 \\
    \bottomrule
  \end{tabular}
\end{table}

% ---------------------------------------------------------------------------
\section{What we are seeing}

\begin{itemize}
  \item \textbf{Normal photos are easy.} The baseline reaches 0.975 on unchanged
        images, matching the $\approx0.98$ reported for this method.
  \item \textbf{A global brightness change costs nothing.} Uniform brightness also
        scores 0.975. LBP compares each pixel only with its neighbours, so an overall
        brightness offset cancels out.
  \item \textbf{One-side lighting breaks it.} As the lighting ramp gets stronger the
        accuracy falls steadily, from 0.960 at $s=0.2$ down to \textbf{0.705} at
        $s=0.8$ and 0.335 at $s=1.0$. A smooth brightness \emph{gradient} across the
        face flips the ``brighter/darker than the centre'' comparison on the darker
        side, so the LBP codes themselves change.
  \item \textbf{Rotation breaks it.} Accuracy falls with the angle, from 0.970 at
        $5^\circ$ to \textbf{0.630} at $20^\circ$ and 0.340 at $30^\circ$. Rotating the
        face rotates the sampling ring around each pixel, which cyclically shifts the
        LBP bit pattern, so the same texture lands in different histogram bins and
        windows.
\end{itemize}

In short, the baseline is accurate on normal photos but fragile to exactly the two
nuisance factors we care about: a lighting gradient and in-plane rotation.

% ---------------------------------------------------------------------------
\section{The experiment today}

Starting from this baseline, the goal is to make the accuracy \emph{stay high} under
the change conditions without losing accuracy on normal photos.

\begin{enumerate}
  \item \textbf{Lighting:} normalise the illumination before the LBP step. Because a
        one-side light is a smooth multiplicative field, dividing the image by its own
        Gaussian-blurred copy (flat-field correction) cancels it while keeping the
        local face texture LBP encodes.
  \item \textbf{Rotation:} make the descriptor tolerant to rotation. One option is a
        rotation-invariant LBP code (take the minimum over the cyclic shifts of the
        code); another is to augment the gallery at enrolment with rotated copies of
        each image so a rotated probe finds a correctly oriented template.
\end{enumerate}

Both variants are evaluated with the \emph{same} split and the \emph{same} conditions
as the baseline above, and compared directly. The experiment succeeds if the rates for
strong one-side lighting and for rotation rise clearly, while the normal-photo rate
stays at its baseline value.

\end{document}
